\documentclass[12pt, letterpaper]{article}

% ── Packages ──────────────────────────────────────────────────────────────────
\usepackage[utf8]{inputenc}
\usepackage[T1]{fontenc}
\usepackage[spanish]{babel}
\usepackage{geometry}
\usepackage{amsmath}
\usepackage{booktabs}
\usepackage{array}
\usepackage{longtable}
\usepackage{graphicx}
\usepackage[breaklinks]{hyperref}
\usepackage{setspace}
\usepackage{titlesec}
\usepackage{fancyhdr}
\usepackage{caption}
\usepackage{float}
\usepackage{parskip}
\usepackage{multirow}
\usepackage{xcolor}
\usepackage{enumitem}
\usepackage{url}

% ── Page geometry ──────────────────────────────────────────────────────────────
\geometry{
  top    = 2.5cm,
  bottom = 2.5cm,
  left   = 3cm,
  right  = 3cm
}

\setlength{\headheight}{14pt}

% ── Header / Footer ───────────────────────────────────────────────────────────
\pagestyle{fancy}
\fancyhf{}
\fancyhead[L]{\small Multiplicación de matrices — HPC}
\fancyhead[R]{\small Universidad Tecnológica de Pereira}
\fancyfoot[C]{\thepage}
\renewcommand{\headrulewidth}{0.4pt}

% ── Section formatting ────────────────────────────────────────────────────────
\titleformat{\section}{\large\bfseries}{}{0em}{}[\titlerule]
\titleformat{\subsection}{\normalsize\bfseries}{}{0em}{}

% ── Hyperlink colours ─────────────────────────────────────────────────────────
\hypersetup{
  colorlinks = true,
  linkcolor  = black,
  urlcolor   = blue,
  citecolor  = black
}

% ── Table column type ─────────────────────────────────────────────────────────
\newcolumntype{C}[1]{>{\centering\arraybackslash}p{#1}}

% ══════════════════════════════════════════════════════════════════════════════
\begin{document}

% ── Title page ────────────────────────────────────────────────────────────────
\begin{titlepage}
  \centering
  \pagenumbering{roman}
  \vspace*{2cm}

  {\LARGE\bfseries Multiplicación de matrices de forma concurrente, con hilos y procesos\par}

  \vspace{2cm}

  {\large\textbf{Presentado por:}\\[0.4em]
  Daniel Toro Soto\\
  Juan Camilo Galvis Agudelo}

  \vspace{1.5cm}

  {\large\textbf{Presentado a:}\\[0.4em]
  Ramiro Andrés Barrios Valencia}

  \vspace{1.5cm}

  {\large HPC: High Performance Computing\\[0.3em]
  Ingeniería de Sistemas y Computación\\[0.3em]
  Universidad Tecnológica de Pereira}

  \vfill

  {\large 2025}
\end{titlepage}
\pagenumbering{arabic}

% ── Table of contents ─────────────────────────────────────────────────────────
\tableofcontents
\newpage

% ══════════════════════════════════════════════════════════════════════════════
\section{Resumen}


% ══════════════════════════════════════════════════════════════════════════════
\section{Introducción}


% ══════════════════════════════════════════════════════════════════════════════
\section{Marco conceptual}

\subsection{High Performance Computing}

\subsection{Multiplicacion de matrices}

\subsection{Matrices transpuestas}

% ══════════════════════════════════════════════════════════════════════════════
\section{Marco Contextual}

\subsection{Características de la máquina}

\subsection{Características del entorno de pruebas}

% ══════════════════════════════════════════════════════════════════════════════
\section{Desarrollo}

% ══════════════════════════════════════════════════════════════════════════════
\section{Pruebas}

\subsection{Algoritmo secuencial}

\subsection{Algoritmo secuencial con optimización de memoria}

\subsection{Algoritmo secuencial con optimización de CPU}

\subsection{Algoritmo con 2 hilos}

\subsection{Algoritmo con 4 hilos}

\subsection{Algoritmo con 8 hilos}

\subsection{Algoritmo con 16 hilos}

\subsection{Algoritmo con 2 procesos}

\subsection{Algoritmo con 4 procesos}

\subsection{Algoritmo con 8 procesos}

\subsection{Algoritmo con 16 procesos}
\newpage

% ══════════════════════════════════════════════════════════════════════════════
\section{Gráficas comparativas}
\newpage

% ══════════════════════════════════════════════════════════════════════════════
\section{Conclusiones}
\newpage

% ══════════════════════════════════════════════════════════════════════════════
\begin{thebibliography}{1}

\bibitem{cortez2004}
Cortéz, A. (2004).
\textit{Teoría de la Complejidad Computacional y Teoría de}…
\textit{Revista de Investigación de Sistemas e Informática}, 1(1), 102--105.
\url{https://revistasinvestigacion.unmsm.edu.pe/index.php/sistem/article/view/3216}

\end{thebibliography}

\end{document}